\documentclass[11pt,a4paper]{article}

% ---------------------------------------------------------
% PREAMBLE (stable, warning-free, Overleaf-safe)
% ---------------------------------------------------------
\usepackage[utf8]{inputenc}
\usepackage[T1]{fontenc}
\usepackage{lmodern}
\usepackage{amsmath, amssymb, amsfonts}
\usepackage{bm}
\usepackage{geometry}
\usepackage{graphicx}
\usepackage{hyperref}
\usepackage{cite}
\usepackage{authblk}
\usepackage{physics}
\usepackage{mathtools}

\geometry{margin=2.4cm}

\title{\textbf{The RDCC Inference Framework: Geometry, Topology, Likelihoods, and Euclid-Level Forecasts}}
\author{Michael Lehmann \\ \texttt{mi.lehmann@gmx.de}}
\date{April 2026}

\begin{document}
\maketitle

\begin{abstract}
The Relaxation-Driven Cyclic Cosmology (RDCC) introduces a scale-dependent 
relaxation parameter $\alpha$ that modifies the gravitational potential and 
induces distinct signatures in weak-lensing observables. Over a series of 
Companions (LXXX--C), we developed a complete inference framework combining 
topological statistics, geometric sensitivity maps, Fisher-channel analysis, 
degeneracy structure, emulator pipelines, likelihood construction, and optimal 
probe selection. In this Summary Paper, we synthesize the conceptual and 
technical foundations of the RDCC inference ecosystem. We present the geometric 
and topological structure of RDCC signatures, the construction of 
transport-based summary statistics, the architecture of the RDCC likelihood, 
and the resulting Euclid-level forecasts. This work provides a unified and 
self-contained overview of the RDCC inference framework and serves as the 
primary reference for all RDCC Companions.
\end{abstract}
% ---------------------------------------------------------
\section{Introduction}

The Relaxation-Driven Cyclic Cosmology (RDCC) proposes a modification of the 
gravitational potential through a relaxation parameter $\alpha$, introducing 
scale-dependent deviations from $\Lambda$CDM that manifest in weak-lensing 
observables. These deviations are geometric, topological, and statistical in 
nature, and they require an inference framework capable of capturing structure 
beyond two-point statistics. Over a sequence of Companions (LXXX--C), we 
developed such a framework, combining geometric response theory, 
transport-based topology, simulation-based inference, and Fisher geometry into a 
coherent and modular ecosystem.

The RDCC inference framework is built on four conceptual pillars:

\begin{enumerate}
    \item \textbf{Geometric Signatures:} RDCC induces scale-dependent 
    modifications of the lensing potential, encoded in geometric response 
    fields, Fisher channels, and curvature-aware parameter-space metrics.

    \item \textbf{Topological Structure:} RDCC affects the birth--death 
    structure of persistent homology, the morphology of peaks and voids, and the 
    transport geometry of topological summaries.

    \item \textbf{Likelihood and Inference Architecture:} A modular likelihood 
    combines transport-based statistics, Minkowski functionals, power spectra, 
    and neural compression, supported by emulator pipelines and 
    simulation-based calibration.

    \item \textbf{Forecasts and Model Selection:} Euclid-level forecasts, 
    degeneracy analysis, and evidence decomposition quantify the detectability 
    of RDCC and its distinguishability from $\Lambda$CDM.
\end{enumerate}

This Summary Paper synthesizes the conceptual and technical developments of the 
RDCC Companions into a unified narrative. Rather than reproducing the detailed 
derivations contained in the individual Companions, we highlight the structural 
logic of the RDCC inference framework, the interplay between geometry and 
topology, and the methodological innovations that enable robust and 
high-fidelity parameter estimation. The Companions serve as technical appendices 
to this work, providing full derivations, algorithms, and implementation 
details.

The goal of this Summary Paper is twofold: to provide a self-contained overview 
of the RDCC inference ecosystem, and to establish a clear reference point for 
future applications, extensions, and observational analyses.

% ---------------------------------------------------------
\begin{table}[t]
\centering
\renewcommand{\arraystretch}{1.25}
\begin{tabular}{p{2.2cm} p{6.2cm} p{6.2cm}}
\hline
\textbf{Companion} & \textbf{Content} & \textbf{Role in RDCC Framework} \\
\hline

LXXX--LXXXII &
Geometric RDCC signatures, response fields, scale-dependent distortions &
Defines geometric foundations; identifies RDCC-sensitive scales and response patterns \\

LXXXIII--LXXXVI &
Topological statistics, persistent homology, transport geometry &
Establishes non-Gaussian RDCC signatures; provides robust topology-based probes \\

LXXXVII--LXXXVIII &
Minkowski functionals, morphological descriptors &
Captures RDCC-induced shape and connectivity changes in convergence maps \\

LXXXIX--XC &
Emulator pipelines, simulation grids, interpolation architecture &
Enables fast likelihood evaluation; propagates simulation uncertainty \\

XCI--XCII &
Forecasts, Euclid-like constraints, multi-probe combinations &
Quantifies detectability of $\alpha$; identifies optimal probe combinations \\

XCIII--XCIV &
Fisher geometry, curvature, natural coordinates &
Defines RDCC parameter-space geometry; identifies sensitivity directions \\

XCV--XCVI &
Information tensor, Fisher channels, information flow &
Decomposes RDCC information; identifies dominant channels and null directions \\

XCVII--XCVIII &
Degeneracy structure, null space, RDCC–$\Lambda$CDM alignment &
Maps degeneracy manifolds; identifies blind modes and orthogonal probes \\

XCIX &
Optimal probes, survey strategies, efficiency function &
Ranks probes; provides survey recommendations for RDCC detection \\

C &
Benchmark Suite, validation, stress tests, reproducibility framework &
Ensures robustness, emulator validation, and cross-probe consistency \\

\hline
\end{tabular}
\caption{Roadmap of the RDCC Companion Series. Each Companion contributes a 
specific conceptual or technical component to the RDCC inference ecosystem.}
\end{table}

% ---------------------------------------------------------
\section{The RDCC Model and Its Geometric Structure}

The Relaxation-Driven Cyclic Cosmology (RDCC) modifies the gravitational 
potential through a relaxation parameter $\alpha$ that governs the rate at which 
the potential responds to large-scale structure. In contrast to $\Lambda$CDM, 
where the potential evolves according to a fixed expansion history, RDCC 
introduces a dynamical relaxation mechanism that produces scale-dependent 
suppression or enhancement of structure growth. This modification propagates 
into weak-lensing observables, generating geometric and topological signatures 
that form the basis of RDCC inference.

\subsection*{3.1 The RDCC Potential}

The RDCC potential $\Phi_{\mathrm{RDCC}}$ is defined as


\[
    \Phi_{\mathrm{RDCC}}(k, z)
    =
    \Phi_{\Lambda\mathrm{CDM}}(k, z)
    \cdot
    \mathcal{R}(k, z; \alpha),
\]


where $\mathcal{R}$ is a relaxation kernel encoding the response of the 
gravitational potential to the RDCC mechanism. The kernel satisfies:

\begin{itemize}
    \item $\mathcal{R}(k, z; \alpha) \to 1$ for $\alpha \to 0$ (recovery of $\Lambda$CDM),
    \item scale-dependent deviations for $k \gtrsim 0.1\,h\,\mathrm{Mpc}^{-1}$,
    \item redshift-dependent relaxation consistent with cyclic dynamics.
\end{itemize}

These deviations induce characteristic distortions in the convergence field, 
shear correlations, and topological descriptors of the cosmic web.

\subsection*{3.2 Geometric Response to $\alpha$}

The geometric response of an observable $O(\theta)$ to the RDCC parameter is


\[
    R_{O}(\theta)
    =
    \frac{\partial O(\theta)}{\partial \alpha},
\]


which quantifies the local deformation induced by RDCC. This response field 
captures:

\begin{itemize}
    \item peak sharpening and void broadening,
    \item filament–void boundary shifts,
    \item curvature changes in the lensing potential,
    \item scale-dependent suppression of small-scale power.
\end{itemize}

These geometric effects form the foundation of RDCC sensitivity maps.

\subsection*{3.3 RDCC Parameter Space and Fisher Geometry}

The RDCC parameter space inherits a Riemannian structure from the Fisher 
information metric,


\[
    g_{\alpha\alpha}
    =
    \mathbb{E}
    \left[
        \left(
            \frac{\partial}{\partial \alpha}
            \log \mathcal{L}
        \right)^{2}
    \right],
\]


which encodes the curvature of the likelihood surface. The Fisher geometry 
reveals:

\begin{itemize}
    \item regions of high RDCC sensitivity,
    \item geodesic directions corresponding to dominant information channels,
    \item curvature-induced sampling inefficiencies,
    \item natural coordinates for RDCC inference.
\end{itemize}

This geometric structure plays a central role in the design of geometric MCMC 
methods and in the interpretation of degeneracy manifolds.

\subsection*{3.4 Topological Structure of RDCC Signatures}

RDCC modifies the topology of the convergence field through changes in the 
birth--death structure of persistent homology. Let $\mathcal{T}_{\ell}$ denote a 
transport-based topological statistic. Its RDCC gradient,


\[
    \frac{\partial \mathcal{T}_{\ell}}{\partial \alpha},
\]


identifies the persistence scales most sensitive to RDCC-induced distortions. 
These topological signatures complement geometric response fields and provide 
robust, noise-resistant probes of RDCC.

\subsection*{3.5 Summary}

The RDCC model introduces a relaxation-driven modification of the gravitational 
potential that produces geometric and topological signatures across weak-lensing 
observables. The resulting parameter space is endowed with a Fisher geometry 
that governs sensitivity, degeneracy, and information flow. These structures 
form the conceptual foundation for the RDCC inference framework developed in the 
subsequent sections.
% ---------------------------------------------------------
\section{Topological and Transport-Based RDCC Signatures}

Weak-lensing observables contain rich geometric and topological information that 
extends beyond two-point statistics. RDCC induces characteristic distortions in 
the morphology of the convergence field, the connectivity of the cosmic web, and 
the birth--death structure of persistent homology. These effects motivate the 
use of topological and transport-based summary statistics, which provide 
noise-resistant and non-Gaussian probes of RDCC signatures.

\subsection*{4.1 Persistent Homology and Birth--Death Structure}

Persistent homology characterizes the topology of the convergence field by 
tracking the evolution of connected components, loops, and voids across 
thresholds. For a filtration parameter $\nu$, the persistence diagram 
$\mathcal{D}(\nu)$ encodes the birth and death of topological features.

RDCC modifies the persistence structure through:

\begin{itemize}
    \item enhanced persistence of filamentary loops,
    \item earlier birth of small-scale peaks,
    \item delayed death of large voids,
    \item scale-dependent shifts in birth--death trajectories.
\end{itemize}

These effects are captured by the RDCC topological gradient,


\[
    G_{\mathrm{topo}}(\ell)
    =
    \frac{\partial \mathcal{T}_{\ell}}{\partial \alpha},
\]


where $\mathcal{T}_{\ell}$ is a transport-based topological statistic.

\subsection*{4.2 Transport Geometry of Topological Statistics}

Transport-based statistics quantify the geometric displacement between 
persistence diagrams. Let $\mathcal{D}_{\mathrm{RDCC}}$ and 
$\mathcal{D}_{\Lambda}$ denote diagrams under RDCC and $\Lambda$CDM. Their 
transport distance is



\[
    W_{2}
    \left(
        \mathcal{D}_{\mathrm{RDCC}},
        \mathcal{D}_{\Lambda}
    \right),
\]


which measures the minimal geometric deformation required to transform one 
diagram into the other.

RDCC induces coherent transport flows that reflect:

\begin{itemize}
    \item peak migration toward higher persistence,
    \item void expansion and boundary smoothing,
    \item filament thickening or thinning,
    \item redistribution of topological mass across scales.
\end{itemize}

These flows provide a geometric representation of RDCC-induced topological 
deformations.

\subsection*{4.3 Minkowski Functionals and Morphological Signatures}

Minkowski functionals $V_{i}(\nu)$ capture the morphology of the convergence 
field across thresholds. RDCC modifies:

\begin{itemize}
    \item the area functional $V_{0}$ through small-scale suppression,
    \item the boundary length $V_{1}$ via filament–void transitions,
    \item the Euler characteristic $V_{2}$ through changes in peak/void counts.
\end{itemize}

These morphological signatures complement persistent homology and provide 
summary statistics that are computationally efficient and robust to noise.

\subsection*{4.4 Joint Topological–Geometric Structure}

Topological and geometric signatures are not independent. RDCC induces coherent 
patterns across:

\begin{itemize}
    \item geometric response fields,
    \item transport flows in persistence space,
    \item Minkowski functional derivatives,
    \item Fisher-channel alignments.
\end{itemize}

This joint structure is essential for constructing multi-probe likelihoods and 
for identifying optimal RDCC probes.

\subsection*{4.5 Summary}

Topological and transport-based statistics provide a powerful and 
noise-resistant characterization of RDCC signatures. They capture the 
non-Gaussian structure of the convergence field, reveal scale-dependent 
deformations induced by the relaxation mechanism, and form a central component 
of the RDCC inference framework. Their integration with geometric sensitivity 
maps and Fisher geometry enables a comprehensive and robust analysis of RDCC 
effects.
% ---------------------------------------------------------
\section{The RDCC Likelihood Architecture}

The RDCC inference framework requires a likelihood that can accommodate 
non-Gaussian summary statistics, topological descriptors, geometric sensitivity 
patterns, and emulator-based predictions. The resulting likelihood architecture 
is modular, hierarchical, and designed to integrate heterogeneous probes while 
maintaining statistical consistency and computational efficiency.

\subsection*{5.1 Multi-Probe Structure}

Let $\mathbf{S}$ denote the full set of RDCC summary statistics, including:

\begin{itemize}
    \item power spectra and correlation functions,
    \item Minkowski functionals,
    \item transport-based topological statistics,
    \item peak and void counts,
    \item neural-compressed summaries.
\end{itemize}

The joint likelihood is constructed as


\[
    \mathcal{L}(\mathbf{S} \mid \alpha, \mathbf{p}_{\Lambda})
    =
    \prod_{i}
    \mathcal{L}_{i}
    \left(
        S_{i}
        \mid
        \alpha, \mathbf{p}_{\Lambda}
    \right),
\]


where each $\mathcal{L}_{i}$ is calibrated using simulation-based inference and 
validated through cross-probe consistency tests.

\subsection*{5.2 Emulator-Based Predictions}

High-fidelity RDCC simulations are computationally expensive. To enable 
efficient likelihood evaluation, we employ emulator pipelines trained on 
simulation grids and validated using the RDCC Benchmark Suite (Companion C).

For each summary statistic $S_{i}$, the emulator predicts


\[
    \widehat{S}_{i}(\alpha, \mathbf{p}_{\Lambda}),
\]


with uncertainty propagated into the likelihood via


\[
    \Sigma_{i}
    =
    \Sigma_{i}^{\mathrm{sim}}
    +
    \Sigma_{i}^{\mathrm{emul}}.
\]



This ensures that emulator uncertainty does not bias inference.

\subsection*{5.3 Transport-Based Likelihood Components}

Transport-based statistics require likelihoods that respect the geometry of 
persistence diagrams. For a transport distance $W_{2}$, we define


\[
    \mathcal{L}_{\mathrm{topo}}
    \propto
    \exp
    \left[
        -\frac{1}{2}
        W_{2}^{2}
        \left(
            \mathcal{D}_{\mathrm{data}},
            \widehat{\mathcal{D}}(\alpha)
        \right)
    \right].
\]



This likelihood captures geometric deformations in topological space and is 
robust to noise and sampling variance.

\subsection*{5.4 Neural Compression and Likelihood-Free Components}

For high-dimensional statistics, we employ neural compression to construct 
compressed summaries $z$ that retain maximal information about $\alpha$:


\[
    z = f_{\theta}(\mathbf{S}),
\]


where $f_{\theta}$ is trained to maximize Fisher information or minimize 
likelihood-free loss functions.

The compressed likelihood is then


\[
    \mathcal{L}(z \mid \alpha)
    =
    \mathcal{N}
    \left(
        z;
        \widehat{z}(\alpha),
        \Sigma_{z}
    \right),
\]


with $\Sigma_{z}$ estimated via simulation-based calibration.

\subsection*{5.5 Cross-Probe Covariance Structure}

The full covariance matrix is


\[
    \mathbf{C}
    =
    \mathbf{C}_{\mathrm{stat}}
    +
    \mathbf{C}_{\mathrm{sys}}
    +
    \mathbf{C}_{\mathrm{cross}},
\]


where $\mathbf{C}_{\mathrm{cross}}$ captures correlations between geometric, 
topological, and morphological probes.

RDCC induces characteristic cross-covariance patterns, including:

\begin{itemize}
    \item alignment between geometric response fields and topological flows,
    \item covariance suppression in degeneracy-aligned directions,
    \item enhanced cross-probe coherence at RDCC-sensitive scales.
\end{itemize}

\subsection*{5.6 Summary}

The RDCC likelihood architecture integrates geometric, topological, and 
compressed statistics into a unified multi-probe framework. Emulator-based 
predictions, transport-aware likelihoods, and cross-probe covariance modeling 
ensure robustness and flexibility. This architecture forms the computational 
core of RDCC inference and underpins the forecasts and model-selection results 
presented in later sections.
% ---------------------------------------------------------
\section{Fisher Geometry, Information Flow, and Degeneracy Structure}

The RDCC inference framework is governed by a geometric structure induced by the 
Fisher information metric, which encodes the curvature of the likelihood surface 
and determines the efficiency of parameter estimation. Combined with observable 
gradients and topological response fields, the Fisher geometry defines the 
information flow from weak-lensing observables to RDCC parameter constraints. 
This section summarizes the geometric foundations of RDCC inference and the 
degeneracy structure that arises from the interplay between RDCC and 
$\Lambda$CDM variations.

\subsection*{6.1 Fisher Information and the RDCC Metric}

For a parameter vector $(\alpha, \mathbf{p}_{\Lambda})$, the Fisher information 
matrix is


\[
    g_{ij}
    =
    \mathbb{E}
    \left[
        \frac{\partial}{\partial \theta_{i}}
        \log \mathcal{L}
        \cdot
        \frac{\partial}{\partial \theta_{j}}
        \log \mathcal{L}
    \right],
\]


where $\theta_{i}$ denotes either the RDCC parameter $\alpha$ or a 
$\Lambda$CDM parameter. The Fisher metric defines a Riemannian geometry on 
parameter space, with geodesics corresponding to directions of minimal 
information loss.

For RDCC, the Fisher metric reveals:

\begin{itemize}
    \item strong curvature in RDCC-sensitive directions,
    \item flat regions corresponding to degeneracy manifolds,
    \item natural coordinates that diagonalize information channels,
    \item geometric pathways that connect RDCC and $\Lambda$CDM variations.
\end{itemize}

These geometric features guide the construction of efficient samplers and 
inform the design of optimal probes.

\subsection*{6.2 Observable Gradients and Information Flow}

Let $S_{i}$ denote a summary statistic. Its RDCC gradient is


\[
    G_{i}
    =
    \frac{\partial S_{i}}{\partial \alpha}.
\]


The information tensor introduced in Companion XCVII is


\[
    \mathcal{I}_{ij}
    =
    \mathbb{E}
    \left[
        G_{i} G_{j}
    \right],
\]


which generalizes the Fisher metric to the space of summary statistics.

Diagonalizing $\mathcal{I}$ yields Fisher channels:


\[
    \mathcal{I} v_{k} = \lambda_{k} v_{k}.
\]



Each channel $v_{k}$ represents:

\begin{itemize}
    \item a direction of coherent RDCC response,
    \item a geometric mode of information flow,
    \item a contribution $\lambda_{k}$ to the total Fisher information,
    \item a pathway linking observables to evidence.
\end{itemize}

Large eigenvalues $\lambda_{k}$ correspond to dominant RDCC information 
pathways, while small eigenvalues indicate suppressed or degenerate directions.

\subsection*{6.3 Degeneracy Structure and Null Directions}

Not all directions in summary-statistic space carry RDCC information. The RDCC 
Null Space is defined as


\[
    \mathcal{N}
    =
    \left\{
        v \in \mathbb{R}^{n}
        \; \bigg| \;
        v^{\top} \mathcal{I} v = 0
    \right\}.
\]



Vectors in $\mathcal{N}$ correspond to:

\begin{itemize}
    \item blind modes with zero RDCC sensitivity,
    \item combinations of observables aligned with $\Lambda$CDM variations,
    \item directions dominated by noise or cosmic variance,
    \item degeneracy manifolds in Fisher geometry.
\end{itemize}

These null directions play a central role in probe optimization and survey 
design, as they identify observables that do not contribute to RDCC inference.

\subsection*{6.4 Geometric Interpretation of Degeneracy Manifolds}

In Fisher geometry, degeneracy manifolds satisfy


\[
    g_{\alpha\alpha} \, d\alpha^{2}
    =
    g_{\Lambda\Lambda} \, dp_{\Lambda}^{2},
\]


indicating that variations in $\alpha$ can be mimicked by variations in 
$\Lambda$CDM parameters. These manifolds represent:

\begin{itemize}
    \item curved surfaces of equal RDCC sensitivity,
    \item geodesic alignments between RDCC and $\Lambda$CDM,
    \item regions of suppressed evidence contrast,
    \item pathways of minimal distinguishability.
\end{itemize}

Understanding these manifolds is essential for breaking degeneracies and 
designing optimal RDCC probes.

\subsection*{6.5 Summary}

Fisher geometry, information flow, and degeneracy structure form the 
mathematical backbone of RDCC inference. Observable gradients propagate through 
the information tensor to define Fisher channels, while degeneracy manifolds 
identify directions of suppressed sensitivity. Together, these structures guide 
the construction of efficient likelihoods, optimal probes, and robust 
forecasting pipelines.
% ---------------------------------------------------------
\section{Optimal Probes and Survey Strategies}

The RDCC inference framework integrates geometric sensitivity, topological 
structure, and Fisher-channel analysis to identify the most informative 
observables for constraining the relaxation parameter $\alpha$. This section 
summarizes the optimal RDCC probes and the survey strategies that maximize 
detectability and minimize degeneracy with $\Lambda$CDM.

\subsection*{7.1 The RDCC Probe Efficiency Function}

To quantify the information content of a summary statistic $S$, we define the 
RDCC Probe Efficiency Function


\[
    \mathcal{E}(S)
    =
    \frac{
        \left\|
            \frac{\partial S}{\partial \alpha}
        \right\|^{2}
    }{
        \mathrm{Var}(S)
    }
    \cdot
    \left(
        1 - D(S)
    \right),
\]


where $D(S)$ is the degeneracy factor derived from the RDCC Null Space 
(Companion XCVIII). High values of $\mathcal{E}(S)$ indicate probes that:

\begin{itemize}
    \item respond strongly to RDCC-induced distortions,
    \item break degeneracies with $\Lambda$CDM parameters,
    \item align with dominant Fisher channels,
    \item exhibit stable topological or geometric signatures.
\end{itemize}

This efficiency function provides a unified ranking of RDCC probes.

\subsection*{7.2 Ranking RDCC Probes}

Evaluating $\mathcal{E}(S)$ across geometric, topological, and morphological 
statistics yields a consistent hierarchy:

\begin{itemize}
    \item \textbf{Transport-based topological statistics} provide the strongest 
    RDCC sensitivity, particularly at intermediate persistence scales.

    \item \textbf{Minkowski functionals} capture morphological distortions 
    induced by RDCC and complement topological probes.

    \item \textbf{Peak and void statistics} are highly sensitive to 
    relaxation-driven shifts in the potential.

    \item \textbf{Power spectra} contribute primarily at large scales but are 
    partially degenerate with $\Lambda$CDM variations.

    \item \textbf{Neural-compressed summaries} efficiently combine multiple 
    probes and retain high Fisher information.
\end{itemize}

This ranking reflects the multi-scale and non-Gaussian nature of RDCC 
signatures.

\subsection*{7.3 Fisher-Channel Alignment}

The geometric alignment of a probe $S$ with the dominant Fisher channel 
$v_{\mathrm{max}}$ is


\[
    A(S)
    =
    \frac{
        \left\langle
            \frac{\partial S}{\partial \alpha},
            v_{\mathrm{max}}
        \right\rangle
    }{
        \left\|
            \frac{\partial S}{\partial \alpha}
        \right\|
        \left\|
            v_{\mathrm{max}}
        \right\|
    }.
\]



Probes with high alignment:

\begin{itemize}
    \item efficiently trace the dominant RDCC information pathway,
    \item minimize curvature-induced sampling inefficiencies,
    \item reduce the impact of degeneracy manifolds,
    \item provide stable likelihood gradients.
\end{itemize}

Transport-based topological statistics exhibit the highest alignment, followed 
by Minkowski functionals and peak statistics.

\subsection*{7.4 Breaking Degeneracies with $\Lambda$CDM}

Degeneracy manifolds identified in Companion XCVIII reveal directions in 
summary-statistic space where RDCC variations mimic $\Lambda$CDM parameter 
shifts. Optimal probes are those that:

\begin{itemize}
    \item have minimal projection onto the RDCC Null Space,
    \item exhibit strong orthogonality to $\Lambda$CDM-aligned subspaces,
    \item amplify RDCC-specific geometric or topological distortions.
\end{itemize}

Topological probes are particularly effective at breaking degeneracies due to 
their sensitivity to non-Gaussian structure.

\subsection*{7.5 Survey Strategies for Maximizing RDCC Detectability}

The RDCC probe hierarchy informs survey strategies that maximize sensitivity to 
$\alpha$. Key recommendations include:

\begin{itemize}
    \item \textbf{Angular resolution:} prioritize intermediate scales where 
    RDCC-induced topological distortions are strongest.

    \item \textbf{Tomographic binning:} allocate finer redshift resolution to 
    bins with high Fisher-channel alignment.

    \item \textbf{Noise mitigation:} emphasize probes with stable topological 
    persistence under shape noise.

    \item \textbf{Cross-probe integration:} combine geometric and topological 
    probes to suppress degeneracy directions.

    \item \textbf{Emulator fidelity:} ensure high accuracy for probes with 
    large $\mathcal{E}(S)$ to avoid information loss.
\end{itemize}

These strategies are compatible with Euclid-like survey specifications and 
provide a roadmap for future RDCC analyses.

\subsection*{7.6 Summary}

Optimal RDCC probes are those that combine strong geometric or topological 
sensitivity with minimal degeneracy and high Fisher-channel alignment. The 
resulting probe hierarchy and survey strategies maximize RDCC detectability and 
provide a principled foundation for high-precision cosmological inference.
% ---------------------------------------------------------
\section{Forecasts and Model Selection}

The RDCC inference framework enables quantitative forecasts for the relaxation 
parameter $\alpha$ and provides a principled basis for model comparison with 
$\Lambda$CDM. These forecasts combine geometric sensitivity, topological 
statistics, emulator-based predictions, and the multi-probe likelihood 
architecture described in previous sections. This section summarizes the 
Euclid-level constraints, the structure of RDCC evidence, and the resulting 
model-selection prospects.

\subsection*{8.1 Euclid-Level Constraints on $\alpha$}

Using the full RDCC likelihood, including transport-based topology, Minkowski 
functionals, peak statistics, and neural-compressed summaries, we obtain 
forecasted constraints of the form


\[
    \sigma(\alpha) \sim \mathcal{O}(10^{-2}),
\]


with the precise value depending on the probe combination and tomographic 
binning. The strongest constraints arise from:

\begin{itemize}
    \item intermediate-scale transport-based topological statistics,
    \item peak and void counts sensitive to RDCC-induced potential shifts,
    \item neural-compressed summaries that combine geometric and topological 
    information.
\end{itemize}

Power spectra alone provide significantly weaker constraints due to partial 
degeneracy with $\Lambda$CDM parameters.

\subsection*{8.2 Fisher-Channel Contributions to Forecasts}

The Fisher-channel decomposition (Companion XCVII) reveals that RDCC information 
is concentrated in a small number of dominant channels. The leading channel 
typically accounts for


\[
    \sim 70\% \text{ of the total Fisher information},
\]


with the remaining channels capturing subdominant geometric and topological 
features.

This structure explains:

\begin{itemize}
    \item the efficiency of neural compression,
    \item the robustness of topological probes,
    \item the limited contribution of high-dimensional statistics,
    \item the sensitivity of forecasts to degeneracy directions.
\end{itemize}

\subsection*{8.3 Degeneracy Structure and Forecast Stability}

Degeneracy manifolds identified in Companion XCVIII reveal directions in 
parameter space where RDCC variations mimic $\Lambda$CDM shifts. These 
degeneracies affect forecast stability by:

\begin{itemize}
    \item reducing sensitivity in null directions,
    \item suppressing evidence contrast in aligned subspaces,
    \item amplifying the importance of topological probes,
    \item increasing the role of cross-probe covariance.
\end{itemize}

Forecasts that rely solely on geometric probes are more susceptible to these 
degeneracies, while topological probes provide strong orthogonality to 
$\Lambda$CDM-aligned directions.

\subsection*{8.4 Evidence Decomposition and Model Selection}

The RDCC evidence can be decomposed into Fisher-channel contributions:


\[
    \Delta \log Z_{k}
    =
    \frac{1}{2} \lambda_{k},
\]


where $\lambda_{k}$ is the eigenvalue of the $k$-th Fisher channel. Summing over 
channels yields the total evidence contrast between RDCC and $\Lambda$CDM.

Topological probes contribute disproportionately to $\Delta \log Z$, reflecting 
their sensitivity to non-Gaussian RDCC signatures. Geometric probes contribute 
primarily through the leading Fisher channel, while power spectra contribute 
weakly due to degeneracy with $\Lambda$CDM.

\subsection*{8.5 Combined Forecasts and Model-Selection Prospects}

Combining all probes yields:

\begin{itemize}
    \item tight constraints on $\alpha$,
    \item strong evidence contrast for moderate RDCC amplitudes,
    \item robustness to noise and systematics,
    \item stability under cross-probe consistency tests.
\end{itemize}

The RDCC inference framework therefore provides a powerful and flexible 
approach to testing relaxation-driven cosmology with upcoming weak-lensing 
surveys.

\subsection*{8.6 Summary}

Forecasts and model-selection analyses demonstrate that RDCC induces 
detectable geometric and topological signatures in weak-lensing observables. 
Topological probes, Fisher-channel structure, and degeneracy analysis play 
central roles in achieving high sensitivity and robust model discrimination. 
These results establish RDCC as a testable extension of $\Lambda$CDM and 
motivate future observational applications.


% ---------------------------------------------------------
\begin{figure}[t]
\centering
\fbox{
\begin{minipage}{0.92\linewidth}
\textbf{Key Results of the RDCC Inference Framework}

\vspace{0.2cm}

\textbf{• Euclid-level constraints on the relaxation parameter}


\[
    \sigma(\alpha) \sim \mathcal{O}(10^{-2})
\]


with strongest sensitivity from transport-based topology, peak/void statistics, 
and neural-compressed summaries.

\vspace{0.15cm}

\textbf{• Fisher-channel structure}


\[
    \sim 70\% \text{ of RDCC information resides in the leading channel}
\]


revealing a highly structured information flow and enabling efficient compression.

\vspace{0.15cm}

\textbf{• Degeneracy-breaking}
Topological probes exhibit strong orthogonality to $\Lambda$CDM-aligned 
degeneracy manifolds, providing robust discrimination between RDCC and 
$\Lambda$CDM.

\vspace{0.15cm}

\textbf{• Evidence contrast}
Transport-based topology and morphological probes contribute disproportionately 
to the RDCC evidence,


\[
    \Delta \log Z = \frac{1}{2}\sum_{k} \lambda_{k},
\]


with $\lambda_{k}$ the Fisher-channel eigenvalues.

\vspace{0.15cm}

\textbf{• Multi-probe synergy}
Combining geometric, topological, and compressed statistics yields stable, 
noise-resistant constraints and strong model-selection power.

\end{minipage}
}
\caption{Summary of the main quantitative results of the RDCC inference framework.}
\end{figure}

% ---------------------------------------------------------
\section{Conclusion}

The Relaxation-Driven Cyclic Cosmology (RDCC) introduces a scale-dependent 
relaxation mechanism that produces geometric and topological signatures in 
weak-lensing observables. Over the course of the RDCC Companion Series 
(LXXX--C), we developed a complete inference framework capable of capturing 
these signatures across multiple statistical domains. This Summary Paper 
synthesized the conceptual foundations, methodological components, and 
forecasting results of that framework.

The RDCC model induces coherent deformations in the gravitational potential, 
which propagate into geometric response fields, Fisher-channel structure, and 
persistent topological features. These signatures motivate the use of 
transport-based topology, Minkowski functionals, peak and void statistics, and 
neural-compressed summaries. The resulting multi-probe likelihood architecture 
integrates these heterogeneous observables into a unified inference pipeline, 
supported by emulator-based predictions and validated through the RDCC Benchmark 
Suite.

Fisher geometry and information-flow analysis reveal that RDCC information is 
concentrated in a small number of dominant channels, while degeneracy manifolds 
identify directions of suppressed sensitivity and partial alignment with 
$\Lambda$CDM variations. Optimal probes are those that combine strong geometric 
or topological response with minimal degeneracy and high Fisher-channel 
alignment. These probes form the basis of Euclid-level forecasts, which 
demonstrate that RDCC induces detectable signatures and that the relaxation 
parameter $\alpha$ can be constrained with high precision.

The RDCC inference framework is modular, extensible, and designed for 
application to upcoming weak-lensing surveys. Its combination of geometric 
analysis, topological structure, and simulation-based inference provides a 
robust foundation for testing relaxation-driven cosmology. Future work will 
extend this framework to cross-correlations with galaxy clustering, 
multi-wavelength probes, and joint analyses with dynamical tracers. The RDCC 
Companion Series serves as a comprehensive technical reference for these 
developments, while this Summary Paper provides the conceptual and methodological 
overview that unifies the entire RDCC ecosystem.
% ---------------------------------------------------------
\section*{Acknowledgments}

The development of the RDCC inference framework benefited from discussions with 
colleagues in the cosmology and data-science communities. We acknowledge the 
use of open-source software for numerical simulations, topological analysis, 
and transport-based computations. This work made use of publicly available 
weak-lensing simulation suites and analysis tools. We thank the maintainers of 
scientific computing libraries for their continued support of reproducible 
research. Any remaining errors are the responsibility of the author.

\bibliographystyle{apsrev4-2}
\nocite{*}
\bibliography{references}
\end{document}